% Сборка: pdflatex lec02.tex
\documentclass[a4paper,12pt]{article}

\usepackage{cmap}
\usepackage[T1,T2A]{fontenc}
\usepackage[utf8]{inputenc}
\usepackage[english,russian]{babel}

\usepackage{amsmath,amssymb,amsthm}
\usepackage[a4paper,margin=2.5cm]{geometry}
\usepackage{enumitem}
\usepackage{listings}

% Чёрно-белое оформление кода: комментарии курсивом, без цвета.
\lstdefinestyle{bw}{
  language=Python,
  basicstyle=\ttfamily\small,
  commentstyle=\itshape,
  numbers=none,
  frame=single,
  columns=fullflexible,
  keepspaces=true,
  showstringspaces=false,
  breaklines=true,
  upquote=true,
}
\lstset{style=bw}

\theoremstyle{definition}
\newtheorem{definition}{Определение}
\newtheorem{example}{Пример}
\theoremstyle{plain}
\newtheorem{theorem}{Теорема}
\newtheorem{lemma}{Лемма}
\theoremstyle{remark}
\newtheorem*{remark}{Замечание}

\renewcommand{\proofname}{Доказательство}
\renewcommand{\qedsymbol}{$\blacksquare$}

\DeclareMathOperator*{\argmax}{arg\,max}
\newcommand{\E}{\mathbb{E}}
\newcommand{\D}{\mathbb{D}}
\renewcommand{\P}{\mathbb{P}}

\begin{document}

\begin{center}
  {\Large\bfseries Теория к лабораторной работе 2\par}
  {\large Метод максимального правдоподобия\par}
\end{center}

В лабораторной работе~1 мы познакомились с точечными оценками параметров распределения и их свойствами (несмещенность, состоятельность). Здесь рассматривается общий способ построения оценок~--- \emph{метод максимального правдоподобия} (ММП), и на примере распределения Пуассона разбираются его применение и свойства получаемой оценки.

\section{Выборка и параметрическая модель}

Пусть $\mathbf{X}=\{X_i\}_{i=1}^n$~--- выборка: набор независимых одинаково распределенных случайных величин с общим распределением из некоторого параметрического семейства. Это семейство задается плотностью (для непрерывных величин) или функцией вероятности (для дискретных) $p(x\mid\theta)$, где $\theta$~--- неизвестный параметр (скалярный или векторный). Требуется по наблюдаемой выборке оценить $\theta$.

\emph{Обозначения.} Прописными буквами $X_i$ обозначаются случайные величины (выборка до проведения эксперимента), строчными $x_i$~--- их реализации, то есть конкретные наблюдённые числа, $\mathbf{x}=\{x_i\}_{i=1}^n$. Математическое ожидание, дисперсия и вероятности берутся от случайных величин $X_i$; функция правдоподобия, график и код работают с числами $x_i$. Эти же обозначения используются в заданиях лабораторной работы.

\emph{Оценкой} параметра называют любую функцию выборки $\hat\theta_n=\hat\theta_n(X_1,\dots,X_n)$. Как функция случайных величин оценка сама случайна; подстановка наблюдённых $x_1,\dots,x_n$ дает число $\hat\theta_n(x_1,\dots,x_n)$. Напомним два основных свойства оценки.

\begin{definition}
Оценка $\hat\theta_n$ называется \emph{несмещенной}, если $\E\hat\theta_n=\theta$ при любом допустимом значении $\theta$.
\end{definition}

\begin{definition}
Оценка $\hat\theta_n$ называется \emph{состоятельной}, если она сходится к $\theta$ по вероятности: для любого $\varepsilon>0$
$$\P\left(|\hat\theta_n-\theta|>\varepsilon\right)\to0,\quad n\to\infty.$$
\end{definition}

\section{Функция правдоподобия и оценка максимального правдоподобия}

Идея метода: из всех возможных значений параметра выбрать то, при котором наблюдаемая выборка «наиболее правдоподобна», то есть имеет наибольшую вероятность (для дискретных величин) или наибольшую плотность (для непрерывных).

Из независимости наблюдений совместная плотность (функция вероятности) выборки равна произведению маргинальных: $\prod_{i=1}^n p(x_i\mid\theta)$. Если рассматривать это выражение как функцию $\theta$ при \emph{фиксированных} наблюдаемых $x_i$, получается функция правдоподобия.

\begin{definition}
\emph{Функцией правдоподобия} реализации $\mathbf{x}=\{x_i\}_{i=1}^n$ случайной выборки $\mathbf{X}=\{X_i\}_{i=1}^n$ из распределения с плотностью (или функцией вероятности) $p(x\mid\theta)$ называется функция параметра $\theta$
$$L(\theta;\mathbf{x})=\prod_{i=1}^n p(x_i\mid\theta).$$
\end{definition}

\begin{definition}
\emph{Оценкой максимального правдоподобия} (ОМП) называется значение параметра, при котором функция правдоподобия максимальна:
$$\hat\theta_n=\argmax_{\theta}L(\theta;\mathbf{x}).$$
\end{definition}

\begin{remark}
Функция правдоподобия $L(\theta;\mathbf{x})$~--- не плотность распределения параметра $\theta$: по $\theta$ она, вообще говоря, не нормирована (интеграл по $\theta$ не равен~1). Это функция параметра при фиксированных данных, а не наоборот.
\end{remark}

\subsection{Логарифм правдоподобия}

На практике максимизируют не саму $L(\theta)$, а ее логарифм.

\begin{definition}
\emph{Логарифмом функции правдоподобия} (log-likelihood) называется
$$\ell(\theta)=\ln L(\theta;\mathbf{x})=\sum_{i=1}^n\ln p(x_i\mid\theta).$$
\end{definition}

Логарифм~--- строго монотонно возрастающая функция, поэтому $L(\theta)$ и $\ell(\theta)=\ln L(\theta)$ достигают максимума в одной и той же точке:
$$\argmax_\theta L(\theta)=\argmax_\theta \ell(\theta).$$
При этом переход к логарифму удобен по двум причинам:
\begin{itemize}
  \item произведение $n$ множителей превращается в сумму, а суммы проще дифференцировать (достаточно линейности производной, не нужно правило дифференцирования произведения);
  \item при больших $n$ произведение $n$ вероятностей (каждая меньше~1) стремится к нулю и может выйти за пределы точности чисел с плавающей запятой (underflow), тогда как сумма логарифмов вычисляется без проблем.
\end{itemize}

\subsection{Как искать ОМП}

Если $\ell(\theta)$ дифференцируема по $\theta$, а точка максимума лежит внутри области допустимых значений параметра, то ОМП находят по схеме:
\begin{enumerate}
  \item записать $L(\theta)$ и $\ell(\theta)$;
  \item найти производную $\ell'(\theta)$ и решить \emph{уравнение правдоподобия} $\ell'(\theta)=0$;
  \item проверить, что найденная стационарная точка~--- максимум. Достаточное условие: $\ell''(\hat\theta)<0$. Если $\ell''(\theta)\le0$ при всех допустимых $\theta$ (функция $\ell$ вогнута), то любая стационарная точка~--- точка глобального максимума;
  \item при необходимости сравнить значение $\ell$ с ее значениями на границе области параметров.
\end{enumerate}

Слагаемые в $\ell(\theta)$, не зависящие от $\theta$, не влияют ни на положение максимума, ни на производную, поэтому при построении графика или вычислении ОМП их можно опускать.

\section{Распределение Пуассона}

\begin{definition}
Случайная величина $\xi$ имеет \emph{распределение Пуассона} $\mathrm{Pois}(\lambda)$ с параметром $\lambda>0$, если она принимает целые неотрицательные значения с вероятностями
$$p(k\mid\lambda)=\P(\xi=k)=\frac{\lambda^k e^{-\lambda}}{k!},\quad k=0,1,2,\dots$$
\end{definition}

Распределение Пуассона описывает число событий, происходящих за фиксированный промежуток времени, если события наступают независимо и с постоянной средней интенсивностью $\lambda$ (например, число звонков в колл-центр за час).

Нам понадобятся математическое ожидание и дисперсия:
$$\E\xi=\lambda,\qquad \D\xi=\lambda.$$

\begin{proof}[Вывод]
Для математического ожидания, сдвигая индекс суммирования $m=k-1$ и используя $\sum_{m\ge0}\lambda^m/m!=e^\lambda$:
$$\E\xi=\sum_{k=0}^\infty k\frac{\lambda^k e^{-\lambda}}{k!}=\lambda e^{-\lambda}\sum_{k=1}^\infty\frac{\lambda^{k-1}}{(k-1)!}=\lambda e^{-\lambda}e^{\lambda}=\lambda.$$
Аналогично для факториального момента:
$$\E\,\xi(\xi-1)=\sum_{k=2}^\infty k(k-1)\frac{\lambda^k e^{-\lambda}}{k!}=\lambda^2 e^{-\lambda}\sum_{k=2}^\infty\frac{\lambda^{k-2}}{(k-2)!}=\lambda^2.$$
Поэтому $\E\xi^2=\E\,\xi(\xi-1)+\E\xi=\lambda^2+\lambda$ и
$\D\xi=\E\xi^2-(\E\xi)^2=\lambda^2+\lambda-\lambda^2=\lambda.$
\end{proof}

\section{Свойства математического ожидания и дисперсии}

Для доказательства свойств оценки понадобятся следующие факты. Пусть $\xi_1,\dots,\xi_n$~--- случайные величины, $a_1,\dots,a_n$~--- константы.
\begin{itemize}
  \item \emph{Линейность математического ожидания} (не требует независимости):
  $$\E\Big(\sum_{i=1}^n a_i\xi_i\Big)=\sum_{i=1}^n a_i\,\E\xi_i.$$
  \item \emph{Дисперсия суммы независимых величин} и однородность дисперсии:
  $$\D\Big(\sum_{i=1}^n \xi_i\Big)=\sum_{i=1}^n \D\xi_i\ \ (\text{при независимости}),\qquad \D(a\xi)=a^2\D\xi.$$
\end{itemize}

\begin{theorem}[Неравенство Чебышёва]
Пусть у случайной величины $\eta$ существует конечная дисперсия. Тогда для любого $\varepsilon>0$
$$\P\left(|\eta-\E\eta|>\varepsilon\right)\le\frac{\D\eta}{\varepsilon^2}.$$
\end{theorem}

\begin{proof}
Пусть $\zeta=(\eta-\E\eta)^2\ge0$. По неравенству Маркова $\P(\zeta\ge a)\le\E\zeta/a$ при $a>0$; при $a=\varepsilon^2$ получаем
$$\P\left(|\eta-\E\eta|>\varepsilon\right)\le\P(\zeta\ge\varepsilon^2)\le\frac{\E(\eta-\E\eta)^2}{\varepsilon^2}=\frac{\D\eta}{\varepsilon^2}.$$
\end{proof}

\begin{remark}
Для несмещенной оценки $\E\hat\theta_n=\theta$, поэтому неравенство Чебышёва принимает вид $\P(|\hat\theta_n-\theta|>\varepsilon)\le \D\hat\theta_n/\varepsilon^2$. Отсюда \emph{достаточное условие состоятельности}: если оценка несмещенная и $\D\hat\theta_n\to0$ при $n\to\infty$, то она состоятельна.
\end{remark}

\begin{theorem}[Закон больших чисел]
Пусть $X_1,X_2,\dots$~--- независимые одинаково распределенные случайные величины с конечным математическим ожиданием $\E X_1=\mu$. Тогда выборочное среднее $\bar X_n=\frac1n\sum_{i=1}^n X_i$ сходится к $\mu$ по вероятности (и почти наверное).
\end{theorem}

Закон больших чисел объясняет два наблюдаемых в работе эффекта. Во-первых, \emph{эмпирическая частота} значения $k$ в выборке, $\hat p_k=\frac1n\#\{i:x_i=k\}$, является выборочным средним независимых индикаторов $\mathbf 1\{X_i=k\}$, у которых $\E\mathbf 1\{X_i=k\}=p(k\mid\lambda)$. Поэтому при росте $n$ эмпирические частоты приближают теоретические вероятности. Во-вторых, выборочное среднее $\bar X_n$ сходится к $\E X_1$ (см.\ далее).

\section{Примеры с решениями}

Ниже разобраны задачи теоретической части; на них опирается практическая часть работы.

\begin{example}
Дайте определение функции правдоподобия и ОМП.
\end{example}

\noindent\textbf{Решение.} Функция правдоподобия реализации $\mathbf{x}=\{x_i\}_{i=1}^n$ случайной выборки $\mathbf{X}=\{X_i\}_{i=1}^n$ из распределения с плотностью (или функцией вероятности) $p(x\mid\theta)$~--- это функция параметра $\theta$ при фиксированной реализации:
$$L(\theta;\mathbf{x})=\prod_{i=1}^n p(x_i\mid\theta).$$
Оценка максимального правдоподобия~--- значение параметра, при котором правдоподобие выборки максимально:
$$\hat\theta_n=\argmax_{\theta}L(\theta;\mathbf{x}).$$

\begin{example}
Почему для нахождения ОМП удобнее максимизировать логарифм функции правдоподобия?
\end{example}

\noindent\textbf{Решение.} Логарифм~--- монотонно возрастающая функция, поэтому точка максимума $L(\theta)$ совпадает с точкой максимума $\ln L(\theta)$. При этом переход к логарифму удобнее:
\begin{itemize}
  \item произведение превращается в сумму: $\ln L(\theta)=\sum_{i=1}^n\ln p(x_i\mid\theta)$, а с суммой проще дифференцировать;
  \item при больших $n$ произведение $n$ вероятностей (каждая $<1$) стремится к нулю и может выйти за пределы точности чисел с плавающей запятой, а сумма логарифмов остается вычислимой.
\end{itemize}

\begin{example}\label{ex:poisson-mle}
Найдите функцию правдоподобия, логарифм правдоподобия и ОМП $\hat\lambda_n$ для реализации $\mathbf{x}=\{x_i\}_{i=1}^n$ случайной выборки $\mathbf{X}=\{X_i\}_{i=1}^n$ из $\mathrm{Pois}(\lambda)$, $\lambda>0$. Убедитесь по знаку второй производной, что найденная точка~--- максимум.
\end{example}

\noindent\textbf{Решение.} Для $\mathrm{Pois}(\lambda)$ функция вероятности $p(x\mid\lambda)=\dfrac{\lambda^x e^{-\lambda}}{x!}$, $x=0,1,2,\dots$ Функция правдоподобия реализации:
$$L(\lambda)=\prod_{i=1}^n\frac{\lambda^{x_i}e^{-\lambda}}{x_i!}=\frac{\lambda^{\sum_i x_i}e^{-n\lambda}}{\prod_i x_i!}.$$
Логарифм правдоподобия:
$$\ell(\lambda)=\ln L(\lambda)=\sum_{i=1}^n\left(x_i\ln\lambda-\lambda-\ln(x_i!)\right)=\Big(\sum_{i=1}^n x_i\Big)\ln\lambda-n\lambda-\sum_{i=1}^n\ln(x_i!).$$
Приравниваем производную по $\lambda$ к нулю:
$$\ell'(\lambda)=\frac{\sum_{i=1}^n x_i}{\lambda}-n=0\quad\Longrightarrow\quad \hat\lambda_n=\frac1n\sum_{i=1}^n x_i=\bar x_n.$$
Вторая производная
$$\ell''(\lambda)=-\frac{\sum_{i=1}^n x_i}{\lambda^2}\le0\quad\text{при }\lambda>0$$
(строго меньше нуля, если не все $x_i$ равны нулю), значит, $\ell$ вогнута, и найденная точка~--- максимум. Итак, ОМП параметра распределения Пуассона~--- выборочное среднее $\hat\lambda_n=\bar x_n$ (как функция случайной выборки $\hat\lambda_n=\bar X_n$). \hfill$\blacksquare$

\begin{remark}
Слагаемое $-\sum_i\ln(x_i!)$ не зависит от $\lambda$. При построении графика $\ell(\lambda)$ его можно опустить: график сдвинется по вертикали, но положение максимума не изменится. Если все $x_i=0$, то $\ell(\lambda)=-n\lambda$ убывает на $(0,\infty)$ и максимум не достигается внутри области $\lambda>0$ (супремум достигается при $\lambda\to0$); для выборок разумного объема из $\mathrm{Pois}(\lambda)$ с не слишком малым $\lambda$ эта ситуация практически невозможна.
\end{remark}

\begin{example}
Докажите, что оценка $\hat\lambda_n$ из примера~\ref{ex:poisson-mle} несмещенная.
\end{example}

\noindent\textbf{Решение.} Элементы выборки независимы и одинаково распределены, $\E X_i=\lambda$ (математическое ожидание $\mathrm{Pois}(\lambda)$ равно $\lambda$). По линейности математического ожидания
$$\E\hat\lambda_n=\E\bar X_n=\frac1n\sum_{i=1}^n\E X_i=\frac1n\cdot n\lambda=\lambda.$$
Оценка несмещенная. \hfill$\blacksquare$

\begin{example}
Докажите, что оценка $\hat\lambda_n$ из примера~\ref{ex:poisson-mle} состоятельная.
\end{example}

\noindent\textbf{Решение.} Для $\mathrm{Pois}(\lambda)$ дисперсия $\D X_i=\lambda$. Из независимости элементов выборки
$$\D\hat\lambda_n=\frac{1}{n^2}\sum_{i=1}^n\D X_i=\frac{\lambda}{n}\to0,\quad n\to\infty.$$
По неравенству Чебышёва (оценка несмещенная, $\E\hat\lambda_n=\lambda$) для любого $\varepsilon>0$
$$\P\left(|\hat\lambda_n-\lambda|>\varepsilon\right)\le\frac{\D\hat\lambda_n}{\varepsilon^2}=\frac{\lambda}{n\varepsilon^2}\to0,\quad n\to\infty.$$
Оценка состоятельна. \hfill$\blacksquare$

\begin{remark}
Состоятельность $\hat\lambda_n=\bar X_n$ следует и непосредственно из закона больших чисел: $\bar X_n\to\E X_1=\lambda$ по вероятности. Кроме того, из формулы $\D\hat\lambda_n=\lambda/n$ видно, что разброс оценки по разным выборкам убывает как $1/\sqrt n$ (стандартное отклонение $\sqrt{\lambda/n}$).
\end{remark}

\section{Практическая часть: что и как считать}

Для практических задач нужны следующие инструменты \lstinline|numpy| и \lstinline|scipy|. Ниже приведены и решения практических задач из ноутбука.

\subsection{Генерация выборки и сравнение частот с функцией вероятности}

Воспроизводимая выборка из $\mathrm{Pois}(\lambda)$ создается генератором \lstinline|rng = np.random.default_rng(seed)| и методом \lstinline|rng.poisson(lam, size)|. Эмпирические частоты получают как \lstinline|np.unique(x, return_counts=True)|, деля счетчики на $n$. Теоретическая функция вероятности~--- \lstinline|scipy.stats.poisson.pmf(k, lam)|. Согласно закону больших чисел частоты должны приближать $p(k\mid\lambda)$.

\begin{example}
Сгенерируйте выборку объема $n=1000$ из $\mathrm{Pois}(4.5)$ и сравните эмпирические частоты с теоретической функцией вероятности.
\end{example}

\noindent\textbf{Решение.}
\begin{lstlisting}
import numpy as np
import matplotlib.pyplot as plt
from scipy import stats

rng = np.random.default_rng(42)

lam = 4.5      # true lambda
n = 1000

x = rng.poisson(lam=lam, size=n)

values, counts = np.unique(x, return_counts=True)
freqs = counts / n

k = np.arange(x.max() + 1)

plt.bar(values, freqs, alpha=0.6, label="empirical frequencies")
plt.plot(k, stats.poisson.pmf(k, lam), "o-", color="C1",
         label=r"theoretical $\mathrm{Pois}(\lambda)$")
plt.xlabel("$k$")
plt.legend()
plt.show()
\end{lstlisting}
Столбики эмпирических частот близки к точкам теоретической функции вероятности.

\subsection{Вычисление ОМП}

По примеру~\ref{ex:poisson-mle} ОМП есть выборочное среднее, то есть \lstinline|np.mean(x)|.

\begin{example}
Реализуйте функцию вычисления ОМП параметра $\lambda$ распределения Пуассона.
\end{example}

\noindent\textbf{Решение.}
\begin{lstlisting}
def hat_lambda_mle(x: np.ndarray) -> float:
    lam = np.mean(x)
    return lam


print(hat_lambda_mle(x), np.mean(x))
assert np.isclose(hat_lambda_mle(x), np.mean(x))
\end{lstlisting}

\subsection{График логарифма правдоподобия}

Чтобы построить график $\ell(\lambda)$, нужно вычислить функцию $\ell$ на сетке значений $\lambda$ при фиксированной выборке. Слагаемое $-\sum_i\ln(x_i!)$ не зависит от $\lambda$, и его можно не учитывать (см.\ замечание к примеру~\ref{ex:poisson-mle}). Сетку нужно брать строго положительной (в $\lambda=0$ логарифм не определен) и содержащей истинное значение и ОМП.

\begin{example}
Постройте график $\ell(\lambda)$ для выборки из предыдущего примера, отметьте вертикальной линией $\hat\lambda_n$ и убедитесь, что оно соответствует максимуму графика.
\end{example}

\noindent\textbf{Решение.}
\begin{lstlisting}
def log_likelihood(lam: float, x: np.ndarray) -> float:
    # the term -sum(log(x_i!)) does not depend on lambda
    return x.sum() * np.log(lam) - len(x) * lam


grid = np.linspace(0.1, 2 * lam, 500)
ll = [log_likelihood(l, x) for l in grid]

lam_hat = hat_lambda_mle(x)

plt.plot(grid, ll, label=r"$\ell(\lambda)$")
plt.axvline(lam_hat, color="black", linestyle="--", label=r"$\hat{\lambda}_n$")
plt.xlabel(r"$\lambda$")
plt.ylabel(r"$\ell(\lambda)$")
plt.legend()
plt.show()

print(lam_hat, grid[np.argmax(ll)])
\end{lstlisting}
График унимодален (вогнут), а максимум на сетке \lstinline|grid[np.argmax(ll)]| совпадает с $\hat\lambda_n$ с точностью до шага сетки.

\subsection{Сходимость оценки при росте $n$}

Чтобы проиллюстрировать состоятельность, для одной выборки считают оценку $\hat\lambda_m$ по первым $m$ элементам для $m=1,\dots,n$ и рисуют ее как функцию $m$. Повторив это для нескольких независимых выборок, получаем несколько траекторий. Согласно примерам выше, $\E\hat\lambda_m=\lambda$ и $\D\hat\lambda_m=\lambda/m$: траектории колеблются вокруг $\lambda$, и их разброс уменьшается как $\sqrt{\lambda/m}$.

\begin{example}
Постройте зависимость $\hat\lambda_m$ от $m$ для $5$ независимых выборок объема $1000$ и отметьте истинное $\lambda$.
\end{example}

\noindent\textbf{Решение.}
\begin{lstlisting}
ns = np.arange(1, n + 1)
K = 5   # number of samples

for _ in range(K):
    xk = rng.poisson(lam=lam, size=n)
    plt.plot(ns, [hat_lambda_mle(xk[:m]) for m in ns], alpha=0.7)

plt.axhline(lam, color="black", linestyle="--", label=r"$\lambda$")
plt.xlabel("$n$")
plt.ylabel(r"$\hat{\lambda}_n$")
plt.legend()
plt.show()
\end{lstlisting}

\subsection{Пример вывода}

Эмпирические частоты хорошо приближают теоретическую функцию вероятности $\mathrm{Pois}(\lambda)$. График логарифма правдоподобия $\ell(\lambda)$ унимодален, а его максимум совпадает со значением $\hat\lambda_n$, вычисленным по формуле $\hat\lambda_n=\bar x_n$ (с точностью до шага сетки). Оценки $\hat\lambda_n$ для разных выборок приближаются к истинному $\lambda$ при росте $n$, а разброс между выборками уменьшается~--- это иллюстрация состоятельности оценки ($\D\hat\lambda_n=\lambda/n\to0$).

\section*{Краткая сводка}

\begin{itemize}
  \item $L(\theta;\mathbf x)=\prod_{i=1}^n p(x_i\mid\theta)$, $\ell(\theta)=\sum_{i=1}^n\ln p(x_i\mid\theta)$, $\hat\theta_n=\argmax_\theta L=\argmax_\theta\ell$.
  \item ОМП ищут из уравнения правдоподобия $\ell'(\theta)=0$ с проверкой знака $\ell''$.
  \item Для $\mathrm{Pois}(\lambda)$: $\E\xi=\D\xi=\lambda$, $\ell(\lambda)=\sum x_i\ln\lambda-n\lambda-\sum\ln(x_i!)$, $\hat\lambda_n=\bar x_n$.
  \item $\hat\lambda_n$ несмещенная ($\E\hat\lambda_n=\lambda$) и состоятельная ($\D\hat\lambda_n=\lambda/n\to0$, неравенство Чебышёва).
  \item Эмпирические частоты сходятся к теоретическим вероятностям по закону больших чисел.
\end{itemize}

\end{document}
